\section{Introduction}
\label{sec:introduction}

A holomorphic vector bundle $E$ on a projective variety is ample when
the tautological quotient line bundle $\OO_{\PP(E)}(1)$ is ample;
here $\PP(E)$ parametrizes one-dimensional quotients
\cite[Proposition~3.2]{Hartshorne1966}.
A smooth Hermitian metric is strictly Griffiths-positive if its Chern
curvature is positive in every pair of nonzero tangent and fiber
directions. Such a metric implies ampleness. Griffiths' conjecture asks
for the converse on smooth complex projective varieties
\cite[Problem~(0.9)]{Griffiths1969}.

Set
\[
 S=\PP^1\times\PP^1,\qquad P=\OO_S(1,1),
\]
where $p_1,p_2:S\to\PP^1$ are the projections, and
$\OO_S(r,s)=p_1^*\OO_{\PP^1}(r)\otimes p_2^*\OO_{\PP^1}(s)$.
Let $f_m:S\to S$ raise the homogeneous coordinates in each factor
to their $m$th powers. Thus $\deg f_m=m^2$ and $f_m^*P\simeq P^m$.

\begin{theorem}\label{thm:main}
There is an algebraic vector bundle $G$ of rank two on $S$ such that,
for every positive integer $m$, the bundle
\[
 E_m=f_m^*G\otimes P
\]
is ample. There is an integer $m_0$ such that $E_m$ admits no smooth
Hermitian metric with strictly Griffiths-positive Chern curvature
for every $m\ge m_0$.
\end{theorem}

The bundle $G$ is given explicitly in Section~\ref{sec:construction}.
The ampleness assertion has a quantitative certificate:
Proposition~\ref{prop:ampleness} proves that $\OO_{\PP(E_m)}(n)$ is
very ample for every even $n$ with $n-1\geq3m$.
The threshold $m_0$ is obtained by compactness and is ineffective.
Thus Griffiths' conjecture fails in rank two on the quadric surface.

\subsection*{Historical context and related results}

Hartshorne developed the algebraic theory of ample vector bundles
through symmetric powers, global generation, and the tautological
line bundle \cite{Hartshorne1966}. Griffiths' original question used
the term ``cohomologically positive''; his definition of ``ample'' in
that article included a first-jet generation condition
\cite[Definitions~(0.2)--(0.3), Problem~(0.9)]{Griffiths1969}.
The conjecture above is the modern formulation of that question.
Its significance is the proposed passage from algebraic positivity
to a Hermitian curvature condition on the bundle itself, which can
then enter differential-geometric vanishing arguments.

For line bundles the algebraic and curvature notions agree
\cite[Proposition~1.15]{Umemura1973}. On complex curves, Hartshorne
characterized ampleness by positivity of the degrees of all nonzero
quotient bundles \cite[Theorem~2.4]{Hartshorne1971}.
Umemura proved the metric equivalence first for elliptic curves and
then for compact Riemann surfaces
\cite{Umemura1972,Umemura1973}; see in particular
\cite[Theorem~2.6]{Umemura1973} and the later work of Campana--Flenner
\cite{CampanaFlenner1990}. These results make a surface the first
possible base dimension for the phenomenon in Theorem~\ref{thm:main}.

Several positive results recover strong curvature after modifying the
bundle. Nakano positivity is a stronger curvature condition than
Griffiths positivity. Berndtsson proved Nakano positivity of $E\otimes\det E$ for
ample $E$ through direct-image curvature
\cite[Theorem~1.3]{Berndtsson2009}. Liu--Sun--Yang established
Nakano and dual-Nakano positivity for determinant-twisted symmetric
powers and for sufficiently large symmetric powers of ample bundles
\cite[Theorems~1.1--1.2 and Corollary~1.3]{LiuSunYang2013}.
The determinant twists and symmetric powers are part of these
conclusions; they do not supply a Griffiths-positive metric on $E$
itself.

Analytic approaches to that remaining question include Demailly's
Hermitian--Yang--Mills system and subsequent curvature functionals
\cite{Demailly2021,Demailly2022}. Pingali identified a limitation of
the initial system and solved a related system for direct sums of
ample line bundles on Riemann surfaces
\cite{Pingali2021,Pingali2023}. Murakami subsequently obtained an
analytic proof of the full curve case through the Demailly system
\cite[Theorem~1.2]{Murakami2026}. These developments concern smooth
Hermitian metrics and clarify the analytic difficulty of recovering
them from ampleness.

The requirement that the metric be Hermitian is essential to the
problem. Kobayashi characterized ampleness of $E$ by negatively curved
strongly pseudoconvex complex Finsler structures on $E^*$
\cite[Theorems~4.1 and~5.1]{Kobayashi1975}. His term ``convex''
means positivity of the complex fiber Hessian of the squared metric
\cite[Section~2]{Kobayashi1975}. Their squared norms are
homogeneous of degree two but need not be Hermitian quadratics in
the fiber variable. More recently, Lempert showed that an equivariant
fiberwise rule applied to a prescribed positive tautological metric
can fail to produce even a semipositive Hermitian metric
\cite[Theorem~2.1]{Lempert2026Construction}. His examples are direct sums of a
positive line bundle and also carry positive Hermitian metrics.
That result identifies an obstruction to the specified construction
rule; the obstruction below concerns every smooth strictly positive
metric on the indicated bundles.

In a preprint posted on September 22, 2026, Du--Xie state the
existence of a rank-two ample bundle on the abelian surface
$C\times C$, where $C$ is the elliptic curve $y^2=x^3-x$, without
a smooth Griffiths-semipositive metric
\cite[Theorem~1.1]{DuXie2026Counterexample}. Their accompanying
preprint gives a first-jet characterization of Griffiths positivity
in terms of an adjoint-Levi obstruction
\cite[Theorem~1.1]{DuXie2026FirstJet}.
The construction here is on the rational quadric and supplies a
specified power-map family, an all-$m$ ampleness certificate, and
the general limiting statement of Proposition~\ref{prop:limit}.
Its metric conclusion is the absence of smooth \emph{strictly}
Griffiths-positive metrics for all sufficiently large $m$.

A subsequent preprint of Liu--Wan, posted on September 29, 2026,
states the existence of ample rank-two bundles on every Hirzebruch
surface, including $\PP^1\times\PP^1$, that are slope-stable with
respect to suitable ample polarizations and admit no smooth
strongly pseudoconvex Finsler metric of semipositive Kobayashi curvature
\cite[Theorem~1.1]{LiuWan2026}. That curvature condition is imposed on
the original bundle; Kobayashi's ampleness characterization above uses
negative curvature on its dual. The result concerns a different
construction and a stronger metric exclusion than the strict
Hermitian conclusion proved here.

\subsection*{Proof strategy}

\paragraph{Compatible graph identifications.}
Choose two projective transformations $A,B$ of $\PP^1$ with order
three and no common fixed point, such that $AB$ is conjugate to
$z\mapsto\lambda z$ with $|\lambda|>1$. Their chosen linear lifts,
and the relative lift $A^{-1}B$, have finite order. The two properties
serve different purposes: finite relative orders allow line maps to
agree at graph intersections, whereas the expanding product will
prevent a common invariant measure.
For $M\in\{I,A,B\}$ put
$\Gamma_M=\{(x,Mx):x\in\PP^1\}$ and
$D=\Gamma_I\cup\Gamma_A\cup\Gamma_B$.
For a suitable common multiple $b$ of the finite lift orders, set
$L_1=\OO_S(b,0)$ and $L_2=\OO_S(0,b)$.
The lifted maps identify $L_1^{-1}|_D$ with $L_2^{-1}|_D$;
duality gives a corresponding identification of $L_1|_D$ and $L_2|_D$.
Define $G$ by the pairs of local sections of $L_1\oplus L_2$ that
agree under this identification on $D$.
This is an elementary transformation in the sense of
\cite[Introduction, (0.1)]{Maruyama1986}.
Section~\ref{sec:construction} checks compatibility at every node,
local freeness, and the two line-bundle quotients of $G$.
Their dual line subbundles in $F=G^\vee$ span away from $D$ and
coincide over $D$ according to the prescribed lifted transformations.

\paragraph{Ampleness from the two quotients.}
The two quotient sequences give two filtrations of symmetric powers.
The product of the two kernel line bundles is positive in both rulings.
Alternating the sequences groups the negative twists into powers of
this product; each remaining quotient loses at most one additional
copy of the graph divisor. After pullback by $f_m$, twisting
$\Sym^n(f_m^*G)$ by $P^{n-1}$ offsets that loss for sufficiently
large even $n$.
The line factors are then globally generated and have
vanishing first cohomology. Section~\ref{sec:ampleness} lifts their
sections through the filtration. A relative Veronese map, the base
embedding supplied by one further factor $P$, and a Segre embedding
then give the stated very-ampleness certificate. This
follows the symmetric-power framework of \cite{Hartshorne1966},
while using the two quotients separately also where their kernels
coincide.

\paragraph{A nonzero quadratic limit.}
Suppose that infinitely many $E_m$ have positive metrics. Their dual
squared norms are plurisubharmonic. Average these norms over $f_m$,
using a local $m$th root of a pulled-back \emph{twisting frame} to
absorb the factor $P$. Each average remains a Hermitian quadratic;
the twisting transition factor tends to one as $m\to\infty$.
After a scalar normalization, the coefficient matrices are locally
uniformly bounded. Proposition~\ref{prop:limit} produces a finite,
nonnegative, nonzero plurisubharmonic function $q$ on $F$, quadratic
on every fiber. Convolution of negative singular Hermitian metrics is
developed
in \cite{BerndtssonPaun2008,Raufi2015,PaunTakayama2018}.
Here noncollapse and the assertion on \emph{every} fiber require
separate arguments: the graph relations lie over a divisor and
cannot be read from an almost-everywhere representative.

\paragraph{Incompatible invariances.}
Restricting $q$ to either line subbundle makes its squared seminorm
independent of the other $\PP^1$ coordinate, by the maximum principle.
The graph relations give a single nonzero squared seminorm on
$\OO_{\PP^1}(-b)$, invariant under the lifts of $A$ and $B$.
Total-space plurisubharmonicity and quadratic fiber dependence imply
subharmonicity of the logarithmic weights. Applying
$i\partial\bar\partial$ to them gives a positive measure whose
nonzero total mass is fixed by the degree
$-b$, even when the seminorm has zeros.
Invariance under $AB$ forces this measure to be supported on its
two fixed points. Both $A$ and $B$ preserve this nonempty support.
Since a permutation of at most two points whose order divides three
is trivial, each would fix every support point, giving a common fixed
point. Section~\ref{sec:obstruction} proves this
contradiction, and Section~\ref{sec:conclusion} assembles the theorem.

\subsection*{Conventions and notation}

For a vector bundle $V$, $\PP(V)$ parametrizes one-dimensional
quotients, and $\OO_{\PP(V)}(1)$ is the tautological quotient line.
The projective line $\PP^1$ used to describe the graph transformations
is also identified with the lines in $\CC^2$; its tautological
subbundle is $\OO_{\PP^1}(-1)$.

We use the curvature convention of \cite{Demailly2021}. At a point,
in local coordinates and a frame orthonormal at that point, write
\[
 \Theta_{E,h}
 =i\sum_{j,k,\alpha,\beta}
 c_{jk\alpha\beta}\,dz_j\wedge d\bar z_k
 \otimes e_\alpha^*\otimes e_\beta.
\]
Strict Griffiths positivity means
\[
 \sum_{j,k,\alpha,\beta}
 c_{jk\alpha\beta}\xi_j\bar\xi_k v_\alpha\bar v_\beta>0
 \quad\text{for all }\xi\ne0,\ v\ne0.
\]
The dual has the opposite Griffiths sign. In particular its squared
norm is plurisubharmonic on the total bundle; this standard criterion
is recalled where it is used. Throughout, ``psh'' means
plurisubharmonic, and
$\ddc u=i u_{z\bar z}\,dz\wedge d\bar z$ on a complex curve.
A Hermitian quadratic form is allowed to be positive semidefinite
when this is stated explicitly.
